\documentclass[11pt]{article}

\usepackage[margin=1in]{geometry}
\usepackage{amsmath,amssymb,amsthm,mathtools}
\usepackage{booktabs}
\usepackage[hidelinks]{hyperref}
\usepackage{microtype}

\newtheorem{theorem}{Theorem}[section]
\newtheorem{corollary}[theorem]{Corollary}
\newtheorem{proposition}[theorem]{Proposition}
\newtheorem{lemma}[theorem]{Lemma}
\theoremstyle{definition}
\newtheorem{definition}[theorem]{Definition}
\newtheorem{problem}[theorem]{Open problem}
\theoremstyle{remark}
\newtheorem{remark}[theorem]{Remark}

\newcommand{\R}{\mathbb R}
\newcommand{\N}{\mathbb N}
\newcommand{\LambdaVM}{\Lambda}
\newcommand{\psd}{\succeq 0}
\newcommand{\pd}{\succ 0}

\title{Withdrawn route note: Square Saturation and Cubic Escape\\
in a Prime-Marked Block--Fock Carrier}
\author{}
\date{July 24, 2026 --- route claim withdrawn}

\begin{document}
\maketitle

\begin{abstract}
We isolate a dilation phase transition in the block--Fock geometry of a finite-aperture
Weil source.  Let
\[
 F(r)=\sum_{j\geq0}r^{4j+1}=\frac{r}{1-r^4},
 \qquad 0<r<1.
\]
The normalized overlap of aligned integer blocks at scales $q$ and $q^a$ is $r^{a-1}$,
where $r=q^{-1/2}$, while their natural radial energies are $F(r)$ and $F(r^a)$.
At the square pair $q,q^2$, the resulting $2\times2$ matrix is indefinite.  We prove
that the von Mangoldt source supplies exactly the missing diagonal capacity, but the repaired
matrix saturates Cauchy--Schwarz and has rank one.  At the cubic pair $q,q^3$, no repair is
needed: the natural matrix is positive definite and
\[
 \det
 \begin{pmatrix}
 F(r)&F(r^2)\\ F(r^2)&F(r^3)
 \end{pmatrix}
 =
 \frac{r^8(1-r^4)}{(1-r^8)^2(1-r^{12})}>0.
\]
Both power maps preserve the von Mangoldt mark.  Hence the rank jump is geometric rather than
an artifact of arithmetic reweighting.

This calculation does not identify a predecessor for finite-aperture Weil positivity.
Multiplicative log-convexity of nonnegative graded partition functions already implies the
underlying positivity mechanism. More decisively, cubic dilation remains inside one prime-power
tower and cannot generate exponent-one prime atoms, transport hard aperture, or determine
first-clock sine signs. Earlier versions incorrectly advertised a triadic Riemann-hypothesis
route. That claim is withdrawn. Result is retained only as a local capacity audit.
\end{abstract}

\section{Introduction}

Weil's criterion converts the Riemann hypothesis into positivity of an explicit quadratic
form~\cite{Weil1952,ConnesConsani2020,Suzuki2026}.  Recent finite-aperture programs seek operator
models in which the prime contribution, the archimedean contribution, and the aperture boundary
are represented on one positive carrier~\cite{CCM2025,ConnesVanSuijlekom2025}.  The difficult
step is not writing the separate scalar identities.  It is proving that their proposed
off-diagonal sewing is compatible with the diagonal capacities required by positivity.

This article studies that compatibility at the first nontrivial power dilations.  The starting
point is the odd block--Fock resolvent
\[
 F(r)=r+r^5+r^9+\cdots=\frac{r}{1-r^4}.                 \tag{1.1}
\]
The exponent class $1\bmod4$ is forced by the finite-aperture Gamma channel.  Normalized aligned
blocks at integer scales $m,n$ have overlap
\[
 c(m,n)=\sqrt{\frac{\min(m,n)}{\max(m,n)}}.             \tag{1.2}
\]
The sewn off-diagonal is therefore $F(c(m,n))$.  Independently, the natural radial vector at
scale $q$ has energy
\[
 D(q)=F(q^{-1/2}).                                      \tag{1.3}
\]

Equations (1.2)--(1.3) are individually positive constructions.  They need not be restrictions
of one positive kernel.  Testing that question on $q,q^a$ gives
\[
 C_a(r)=
 \begin{pmatrix}
 F(r)&F(r^{a-1})\\
 F(r^{a-1})&F(r^a)
 \end{pmatrix},
 \qquad r=q^{-1/2}.                                    \tag{1.4}
\]
Our main result is the sharp change between $a=2$ and $a=3$.

\begin{theorem}[Square saturation and cubic escape]
\label{thm:main}
Let $0<r<1$.
\begin{enumerate}
\item The natural square carrier
\[
 C_2(r)=
 \begin{pmatrix}F(r)&F(r)\\F(r)&F(r^2)\end{pmatrix}
\]
is indefinite.
\item Adding $\delta$ to its second diagonal makes it positive semidefinite exactly when
\[
 \delta\geq F(r)-F(r^2).                               \tag{1.5}
\]
At equality, the repaired matrix has rank one.
\item The natural cubic carrier
\[
 C_3(r)=
 \begin{pmatrix}F(r)&F(r^2)\\F(r^2)&F(r^3)\end{pmatrix}
\]
is positive definite, with
\[
 \det C_3(r)=
 \frac{r^8(1-r^4)}{(1-r^8)^2(1-r^{12})}.              \tag{1.6}
\]
\end{enumerate}
\end{theorem}

For an integer source atom, squaring and cubing preserve the von Mangoldt coefficient:
\[
 \LambdaVM(q^2)=\LambdaVM(q)=\LambdaVM(q^3).           \tag{1.7}
\]
We prove below that the square counterterm in (1.5) is not fitted.  It decomposes canonically
into two positive source losses.  Their sum is exactly the threshold, so source repair closes the
square positivity gap but cannot create rank.

The cubic determinant (1.6) is a strict common-carrier surplus in this local power hierarchy.
It does not justify replacing binary completion by triadic valuation-residue completion.
A rank-one repaired square block has zero Schur innovation and positive cubic determinant has a
nonzero local one, but neither statement connects that innovation to full arithmetic source.
Section~\ref{sec:rhroute} records precise obstruction and withdrawal of former route.

\section{Block overlaps and radial capacities}

Let $e_q$ be the normalized indicator of an aligned block of length $q$.  Direct intersection
counting gives
\[
 \langle e_m,e_n\rangle
 =\frac{\min(m,n)}{\sqrt{mn}}
 =c(m,n).                                               \tag{2.1}
\]
Passing to tensor degrees $1,5,9,\ldots$ gives the complete overlap $F(c(m,n))$.
For $n=q^a$ and $r=q^{-1/2}$,
\[
 c(q,q^a)=q^{(1-a)/2}=r^{a-1}.                         \tag{2.2}
\]
The natural radial energy at the second scale is
\[
 D(q^a)=F((q^a)^{-1/2})=F(r^a).                        \tag{2.3}
\]
Equations (2.2)--(2.3) prove (1.4).

\begin{lemma}[Strict radial monotonicity]
\label{lem:mono}
$F$ is strictly increasing on $(0,1)$.
\end{lemma}

\begin{proof}
Differentiation of (1.1) gives
\[
 F'(r)=\frac{1+3r^4}{(1-r^4)^2}>0.
\]
\end{proof}

\section{The square capacity law}

At $a=2$, equations (2.2)--(2.3) give
\[
 C_2(r)=
 \begin{pmatrix}
 F(r)&F(r)\\
 F(r)&F(r^2)
 \end{pmatrix}.                                       \tag{3.1}
\]
Since $0<r^2<r<1$, Lemma~\ref{lem:mono} gives $F(r^2)<F(r)$.  Hence
\[
 (1,-1)C_2(r)(1,-1)^{\mathsf T}=F(r^2)-F(r)<0,         \tag{3.2}
\]
which proves Theorem~\ref{thm:main}(1).

Consider the repaired matrix
\[
 C_2(r;\delta)=
 \begin{pmatrix}
 F(r)&F(r)\\
 F(r)&F(r^2)+\delta
 \end{pmatrix}.                                       \tag{3.3}
\]
Because $F(r)>0$, the $2\times2$ principal-minor criterion yields
\begin{align}
 C_2(r;\delta)\psd
 &\Longleftrightarrow
 F(r)\bigl(F(r^2)+\delta\bigr)-F(r)^2\geq0 \notag\\
 &\Longleftrightarrow
 \delta\geq F(r)-F(r^2).                               \tag{3.4}
\end{align}
At equality,
\[
 C_2\bigl(r;F(r)-F(r^2)\bigr)
 =
 F(r)\begin{pmatrix}1&1\\1&1\end{pmatrix},             \tag{3.5}
\]
so the repaired carrier is positive semidefinite of rank one.  More generally,
\[
 \det C_2(r;\delta)
 =F(r)\bigl(\delta-[F(r)-F(r^2)]\bigr).                \tag{3.6}
\]
Strict two-dimensional positivity is therefore equivalent to strict surplus above the
square-capacity debt.

\section{The von Mangoldt source closes the debt exactly}

Write
\[
 T(r)=F(r)-r=\frac{r^5}{1-r^4}                         \tag{4.1}
\]
for the higher odd block--Fock tail, and
\[
 w(q)=\LambdaVM(q)q^{-1/2}.                            \tag{4.2}
\]
The identity $\LambdaVM(q^2)=\LambdaVM(q)$ holds both on and off prime-power support.

\begin{theorem}[Source-fracture law]
\label{thm:fracture}
For every positive integer $q$, with $r=q^{-1/2}$,
\[
 \LambdaVM(q)\,[F(r)-F(r^2)]
 =
 [w(q)-w(q^2)]
 +\LambdaVM(q)\,[T(r)-T(r^2)].                         \tag{4.3}
\]
For $q\geq2$, both terms on the right are nonnegative.  If
$\LambdaVM(q)\neq0$, their sum divided by $\LambdaVM(q)$ is exactly the minimum counterterm
in \emph{(3.4)}.
\end{theorem}

\begin{proof}
Since $F(r)=r+T(r)$ and
\[
 w(q^2)=\LambdaVM(q^2)(q^2)^{-1/2}
       =\LambdaVM(q)r^2,
\]
expanding the right side of (4.3) gives its left side.  For $q\geq2$,
$0<r^2<r<1$.  Both $r$ and $T(r)=r^5/(1-r^4)$ are increasing on $(0,1)$, while
$\LambdaVM(q)\geq0$.  This proves nonnegativity.  Division by a nonzero von Mangoldt mark and
(3.4) prove the final assertion.
\end{proof}

\begin{corollary}[Canonical square saturation]
\label{cor:saturation}
On every von Mangoldt source atom, the canonical critical-plus-tail square innovation makes the
common carrier positive semidefinite, but its determinant is zero.
\end{corollary}

This is the decisive negative conclusion for the square route.  The missing capacity was real,
and the arithmetic source does supply it, but only at the Cauchy--Schwarz boundary.  No
two-dimensional Schur complement can be extracted from (3.5).

Iterated squaring also obeys a conservation law.  If $q_0=q$ and $q_{j+1}=q_j^2$, then
\[
 \sum_{j=0}^{N-1}[w(q_j)-w(q_{j+1})]
 =w(q)-w(q_N).                                         \tag{4.4}
\]
Thus square innovation is transported down one marked orbit without creating additional mass.
Equation (3.5) shows the parallel absence of geometric rank creation.

\section{Cubic escape}

At $a=3$, the common carrier is
\[
 C_3(r)=
 \begin{pmatrix}
 F(r)&F(r^2)\\
 F(r^2)&F(r^3)
 \end{pmatrix}.                                       \tag{5.1}
\]
Substitution of $F(s)=s/(1-s^4)$ gives
\begin{align}
 \det C_3(r)
 &=
 \frac{r^4}{(1-r^4)(1-r^{12})}
 -\frac{r^4}{(1-r^8)^2}.                              \tag{5.2}
\end{align}
Put $x=r^4$.  Clearing the positive denominators in (5.2) exposes the factorization
\[
 (1-x^2)^2-(1-x)(1-x^3)=x(1-x)^2.                     \tag{5.3}
\]
Cancelling one factor $1-x$ proves (1.6).  Since every factor in (1.6) is positive for
$0<r<1$, the determinant is positive.  Its first diagonal $F(r)$ is also positive, so
$C_3(r)\pd$.

The corresponding cubic Schur innovation is
\[
 \mathcal I_3(r)
 :=F(r^3)-\frac{F(r^2)^2}{F(r)}
 =\frac{\det C_3(r)}{F(r)}>0.                          \tag{5.4}
\]
No counterterm enters (5.4).  Moreover
\[
 \LambdaVM(q^3)=\LambdaVM(q),                          \tag{5.5}
\]
so the strict innovation is compatible with the same source mark as the square law.

\begin{corollary}[Source-marked rank transition]
For every integer $q\geq2$ with $\LambdaVM(q)\neq0$, the canonical source-repaired carrier at
$q,q^2$ has rank one, whereas the natural carrier at $q,q^3$ is positive definite.
\end{corollary}

Its three entries are fixed by aligned block overlap and natural odd-Fock radial energy.
Equality at square threshold and strict factor $r^8(1-r^4)$ at cube are useful normalization
checks, but underlying positivity is standard multiplicative log-convexity.

\section{Why cubic dilation is not a first-clock predecessor}
\label{sec:rhroute}

For logarithmic aperture $L>0$, the first physical sine clock is
\[
 \xi_L=\frac{2\pi}{L}.                                 \tag{6.1}
\]
The finite prime part of the source is
\[
 P_L=
 \sum_{2\leq q\leq e^L}
 \LambdaVM(q)q^{-1/2}\sin(\xi_L\log q).                \tag{6.2}
\]
Midpoint reflection decomposes the complete source coefficient as
\[
 \pi B_1(L)=A_L+P_L,                                   \tag{6.3}
\]
where $A_L>0$ is the continuous Gamma--polar reflection budget.  Equivalently,
\[
 B_1(L)\geq0
 \quad\Longleftrightarrow\quad
 P_L^{-}\leq A_L+P_L^{+},                              \tag{6.4}
\]
with $P_L^\pm$ the two nonnegative Jordan half-wave masses.  Positivity of
$B_1(L)$ for every physical aperture in the required range is an exact
Riemann-hypothesis criterion.

The square carrier does not advance (6.4). Corollary~\ref{cor:saturation} says that its
source-derived Schur innovation vanishes identically. Cubic dilation supplies local strictness:
\[
 \sin(\xi_L\log(q^3))=\sin(3\xi_L\log q).              \tag{6.5}
\]
Prime-power valuations split formally into three affine classes
\[
 3e,\qquad 3e+1,\qquad 3e+2.                           \tag{6.6}
\]
The local carrier (5.1) controls only homogeneous map $p^e\mapsto p^{3e}$ inside each fixed-prime
tower. This loses exactly data required by (6.2):
\begin{enumerate}
\item exponent-one atoms $p$ are not images of positive integral exponents under $e\mapsto3e$;
\item affine sectors $3e+1$ and $3e+2$ require independent initial data rather than consequences
of homogeneous cubic carrier;
\item hard cutoff $p^e\leq e^L$ is not transported by fixed towerwise Gram inequality; and
\item positivity of radial kernel contains no information determining sign of
$\sin(\xi_L e\log p)$.
\end{enumerate}
Therefore no coefficientwise implication from (5.1) to (6.4) has been established. Adding
off-sector maps with required atom, cutoff, and sign data would insert missing first-clock
source theorem rather than derive it. Former ``triadic assembly problem'' is consequently
invalid as a backward RH interface and is withdrawn.

\section{Relation to existing operator programs}

Semigroup Toeplitz and $ax+b$ operator algebras provide general dilation and covariance
frameworks~\cite{BrownloweEtAl2010,LacaSehnem2022}.  They do not determine the diagonal
normalization (1.3), the odd-Fock resolvent (1.1), or the source mark (1.7).  The present rank
transition occurs only when those data are imposed simultaneously.

Connes--Consani and Connes--Consani--Moscovici place finite-aperture Weil forms inside scaling and
spectral-triple constructions~\cite{ConnesConsani2020,CCM2025}.  Connes--van Suijlekom establish
real-zero consequences for suitable lower-bounded finite-interval convolution
operators~\cite{ConnesVanSuijlekom2025}.  Suzuki recasts Weil positivity through screw functions
and nonlocal finite-interval operators~\cite{Suzuki2026}. These works motivate search for a
single positive source carrier.

Algebraic identity (5.3) is an elementary specialization of multiplicative log-convexity for
nonnegative graded partitions. Function $F(r)=\sum_{j\geq0}r^{4j+1}$ is such a partition.
Square threshold is standard $2\times2$ positive-semidefinite completion, while source fracture
uses only $F(r)=r+T(r)$ and $\Lambda(q^2)=\Lambda(q)$. Exact normalization audit is useful, but
does not support scientific priority claim.

\section{Conclusion}

First two source-mark-preserving power dilations lie on opposite sides of local carrier-rank
boundary. Squaring requires exact source debt and stops at rank one; cubing requires no repair and
produces determinant (1.6). This rejects square-chain Schur innovation but does not promote cubic
valuation geometry to RH predecessor.

Former triadic route confused local radial log-convexity with coefficientwise assembly of actual
arithmetic source. Since latter contains independent prime atoms, hard cutoff, and oscillatory
signs absent from cubic carrier, route claim is withdrawn. Rightmost problem remains construction
of a source-complete carrier whose compression is proved, not postulated, to equal first-clock
dominance.

\begin{thebibliography}{99}

\bibitem{BrownloweEtAl2010}
N.~Brownlowe, A.~an Huef, M.~Laca, and I.~Raeburn,
\emph{Boundary quotients of the Toeplitz algebra of the affine semigroup over the natural
numbers},
Ergodic Theory Dynam.\ Systems \textbf{32} (2012), 35--62;
\href{https://arxiv.org/abs/1009.3678}{arXiv:1009.3678}.

\bibitem{CCM2025}
A.~Connes, C.~Consani, and H.~Moscovici,
\emph{Zeta spectral triples},
\href{https://arxiv.org/abs/2511.22755}{arXiv:2511.22755}, 2025.

\bibitem{ConnesConsani2020}
A.~Connes and C.~Consani,
\emph{The scaling Hamiltonian},
J.\ Operator Theory \textbf{85} (2021), 103--131;
\href{https://arxiv.org/abs/1910.14368}{arXiv:1910.14368}.

\bibitem{ConnesVanSuijlekom2025}
A.~Connes and W.~D.~van Suijlekom,
\emph{Quadratic forms, real zeros and echoes of the spectral action},
\href{https://arxiv.org/abs/2511.23257}{arXiv:2511.23257}, 2025.

\bibitem{LacaSehnem2022}
M.~Laca and C.~F.~Sehnem,
\emph{Toeplitz algebras of semigroups},
J.\ Funct.\ Anal.\ \textbf{283} (2022), 109429;
\href{https://arxiv.org/abs/2101.06822}{arXiv:2101.06822}.

\bibitem{Suzuki2026}
M.~Suzuki,
\emph{Weil's quadratic form via the screw function},
\href{https://arxiv.org/abs/2606.09096}{arXiv:2606.09096}, 2026.

\bibitem{Weil1952}
A.~Weil,
\emph{Sur les ``formules explicites'' de la th\'eorie des nombres premiers},
Comm.\ S\'em.\ Math.\ Univ.\ Lund (1952), 252--265.

\end{thebibliography}

\end{document}
